% sec-09: risks, stop conditions, and the cost of the conversion.  No figure.
% FINAL stage: the risk register's mitigation column is re-cut so the three
% rows reading none today or partial sort to the head of the reader's eye, and
% the closing subsection is tightened by one sentence.
\section{Risks, Stop Conditions, and What This Is Not}

\subsection{What the Conversion Costs}

Four things are worse about this document than about the ten it replaces. Each
is stated first and answered second, in that order.

\textbf{1. The individual-scientist framing is abandoned.} It is what made all
ten prior applications distinctive, and it is the framing the report rewards
most explicitly. The answer is that the framing was never the subject; the trial
was. An entity applying under a person-based mechanism would be applying under
the wrong clause, and \S1 shows the eligibility test it fails.

\textbf{2. Capitalization detail is published that the author might prefer to
keep private.} Table 13 states the founder's holding at each of five stages,
Table 9 and Table 10 state every line of both budgets, and all of it is
permanent under a DOI. The answer is that a milestone schedule with a redacted
budget is not auditable, and an unauditable schedule is the thing this document
exists to replace.

\textbf{3. Three of the ten prior recipients will read this as off-mechanism.}
The NIH Director's Pioneer Award and the HHMI Investigator Program are
person-based, and the focused research organization route is
organization-based and time-bound to dissolution. A company plan is the wrong
document for all three. The answer is that it was always the wrong document for
all three, and \S1 is the argument that it is the right one for the fourth.

\textbf{4. A financing paper is not a scientific publication and fits no
journal.} It has no hypothesis, no method, and no result in the sense a journal
uses those words. The answer is that Zenodo is the correct venue for it and the
only one claimed, and that the paper is labelled a plan on its cover.

\subsection{The Risk Register}

Table 19 carries ten risks against the halt point or milestone each attaches to,
and states whether the mitigation exists today or is merely planned.

\begin{apptable}
\begin{tabularx}{\textwidth}{>{\raggedright\arraybackslash}p{3.4cm} >{\raggedright\arraybackslash}p{1.5cm} >{\raggedright\arraybackslash}p{1.9cm} >{\raggedright\arraybackslash}X}
\headrow \thead{Risk} & \thead{Attaches to} & \thead{Likelihood} & \thead{Mitigation, and whether it exists today}\\
\midrule
Site agreement does not execute & H4, G4, M1 & High & None today. Requires an executed CTA that no company action produces \\
IND clinical hold issued & H1, G3, M5 & Low & The IND package is drafted and deposited; the 30-day clock is the only unknown \\
Two of three DLTs at dose level 1 & H2, M8 & Moderate & Exists: 3+3 de-escalation is written into the protocol with a defined dose range \\
Phase II award is not made & H3, M6 & Moderate & Exists: the Descoped state resubmits at three participants rather than six \\
Founder falls below 50 percent & H5, stage 4 & Low inside 33 months & Exists: every tier 3 instrument is gated behind a milestone, and stage 4 sits outside the plan \\
Investigator acquires a \S54.2 interest & H5 & Low & Exists: the investigator role is assigned wholly to the site and each files Form 3454 \\
Advisory output mistaken for a directive & M7 & Moderate & Exists: the model is confined to advisory output and every arrow reaching the robot leaves a human first \\
Bench p95 latency above 250 ms & H1, G1, M3 & Low & Exists: the interlock rig is a Phase I deliverable measured before any participant \\
The 81.9 credibility score fails to reproduce & M4, WP15 & Low & Partial: the suite is hashed at M4, but the full re-run is WP15 and is unfunded at \$1,606,000 \\
No alternate San Diego site accepts & H4 & Moderate & None today. The plan names one site and holds no fallback agreement \\
\end{tabularx}
\end{apptable}
\vspace{-0.65cm}
\tabcap{Table 19. Ten risks, the halt point or milestone each attaches to, and\\
whether the mitigation exists today or is merely planned. Three rows read\\
none today or partial, and those three are the honest content of the table.}

Three rows carry no mitigation that exists today. Two of the three are the same
fact seen twice: the plan depends on one institution and holds no alternative.
The third, the verification re-run, is unfunded because it sits in the
\$2,104,000 the SBIR route does not reach, which is the kind of consequence \S3
exists to make visible rather than to bury.

\subsection{What This Document Is Not}

It is not a submission of record. It is not an agreement with UC San Diego,
Moores Cancer Center, or any other institution, and no agreement of any kind
exists. It is not a completed raise: no term sheet, simple agreement for future
equity, or subscription agreement exists, and every capitalization figure in
Table 13 is a plan. It is not a claim that daraxonrasib is first in human; the
agent is investigational and already in Phase 3 evaluation. It is not medical or
regulatory advice. No robotic configuration has been specified or cleared. No
patient has been treated.

\subsection{The Two Single Points of Failure}

Figure 15 makes them visible, and a reader who has skipped the figure should
still meet them here.

H1, the month-nine gate, is an OR branch. It fails on either a bench stop
latency above 250 ms at the 95th percentile or an IND clinical hold. Either
alone ends Phase I without a Phase II request.

H5, the ownership test, is the other OR branch. It fails on either the founder
falling below 50 percent or a clinical investigator acquiring an interest under
21 CFR \S54.2. Either alone breaks the eligibility the whole plan rests on.

Neither is technical in the sense that more engineering fixes it. One is
regulatory and instrumental; the other is a fact about who owns the shares. Both
are knowable in advance and both are measured before they can do damage, which
is the only reason a schedule carrying two single points of failure is worth
putting in front of a program officer.
